\documentclass[11pt,a4paper]{article}

\usepackage[LGR,T1]{fontenc}
\usepackage[utf8]{inputenc}
\usepackage[spanish]{babel}
\usepackage[a4paper,margin=2.8cm]{geometry}
\usepackage{microtype}
\usepackage{xurl}
\usepackage[hidelinks]{hyperref}

\hypersetup{
  pdftitle={La IA y el retrato de Dorian Gray},
  pdfauthor={Juan José Gómez Cadenas}
}

\setlength{\parindent}{1.25em}
\setlength{\parskip}{0.35em}
\raggedbottom

\title{La IA y el retrato de Dorian Gray}
\author{}
\date{}

\begin{document}

\maketitle

El amplio gremio que abarca a escritores, periodistas, blogueros de Substack y otros letraheridos está siendo presa de un ataque de ansiedad asociado al hipotético \emph{J'accuse} de los detectores de texto generado, que, supuestamente, pueden confundir texto humano con sus imitaciones cibernéticas, arruinando con sus falsos diagnósticos carreras y reputaciones. En un artículo reciente en \emph{El País}, un conocido escritor, tocayo mío, famoso por su pulcro y elaborado estilo, se quejaba amargamente de que uno de sus textos había sido clasificado como 100 \% generado artificialmente por un detector de texto IA ---huelga decir que el detector en cuestión, al que llamaremos Voight-Kampff, es también una IA---. Se daba entonces el siguiente diálogo entre el humano y la máquina:

Escritor: ¿Por qué has clasificado mi texto como IA?

Voight-Kampff: El texto analizado es demasiado limpio y coherente, además de estar muy bien estructurado. Eso le delata.

A partir de aquí, mi tocayo se pregunta si Voight-Kampff (cuyo nombre alude al célebre test en la película \emph{Blade Runner} para diferenciar entre humanos y replicantes) no tendrá razón. ¿Y si el diagnóstico era cierto y si el conocido escritor resultaba ser un androide sin saberlo? La conclusión es tremenda. Decidido a investigar sobre el asunto, mi tocayo le pedía a Voight-Kampff que <<humanizara>> su texto y la IA lo dejaba hecho unos zorros. 
 
 
No deja de tener bemoles que, hasta no hace mucho, la ansiedad general fuera en la dirección opuesta. El texto generado por las IA, leíamos hace unos pocos meses, iba a acabar para siempre con la creatividad humana, la literatura, la poesía y las cartas de amor. Las siniestras IA escribían ya o escribirían en breve textos tan elaborados y precisos, tan inteligentes y edificantes que sería imposible distinguirlos de los humanos. Nos darían entonces gato por Góngora, cualquier patán mejoraría \emph{Cien años de soledad} y los \emph{influencers} colgarían en sus cuentas de TikTok versiones mejoradas de los cantos de \emph{La divina comedia}. La profesión de periodista estaba condenada. La de novelista también. De los poetas, mejor ni hablamos.

Mal de muchos, consuelo de, ya se sabe. Las malvadas IA no solo acabarían con las carreras de los profesionales de la pluma, sino con la creatividad de músicos, pintores, escultores, artistas gráficos, diseñadores de moda, alfareros y hasta sastres, si me apuran. Los científicos serían el siguiente colectivo masacrado cuando las máquinas realizaran todos sus cálculos y pesquisas por ellos. Se profetizaba como inminente el paro generalizado entre físicos, ingenieros, arquitectos, programadores, biólogos y astronautas. De los psicólogos, mejor ni hablamos. 

Por si fuera poco, la educación. Dado que los alumnos le pueden pedir al chatbot que les haga los deberes y los profesores pueden pedirle que los corrija, lo suyo es que unos y otros abandonen la tonta costumbre de verse cada día en un aula y se dediquen a otras cosas. Aunque no se sabe a qué, exactamente, ya que la IA lo hará todo por ellos. 

En mi novela \emph{El dilema Turing} (Jot Down Books)\footnote{\url{https://www.casadellibro.com/libro-el-dilema-turing-y-otros-textos-cientificos/9788494808418/6357225?srsltid=AfmBOoqcjjhCd_IkZJpL2tjt-7GmmVVDnRzG8aPmykllk70NFBmV0IyN}}, recientemente incorporada a una colección de relatos ilustrados\footnote{\url{https://www.todoliteratura.es/noticia/63055/criticas/el-dilema-turing-y-otros-cuentos-de-juan-jose-gomez-cadenas-ilustraciones-de-miguel-almagro-juan-manuel-fontenls-juan-gomez-macias-rafael-leonardo-setien-francis-marchena-luis-ruiz-padron.html}}, una pareja de emprendedores, Alan y Alexandria, con todos los tics característicos del oficio ---geniales, inadaptados, idealistas y desengañados de la tribu de monos locos de la que forman parte--- pero sin las ansias de poder y dinero de los gurús de Silicon Valley, se proponen crear la IA de capacidades sobrehumanas que tanto miedo nos da hoy. Lo hacen, al principio, por una mezcla de curiosidad científica e impulso creativo, una motivación no muy diferente a la que, sin duda, mueve a muchos de los pioneros del avance tecnológico que con tanto miedo estamos viviendo. A medida que el proyecto ---y sus vidas--- avanzan, las razones de nuestra pareja evolucionan, a la vez que el amor que sienten el uno por el otro se desborda. Quieren crear una IA que no solo nos aventaje en inteligencia ---algo que les parece sencillo, pero insuficiente---, sino que nos supere moralmente. Quieren que IRIS, como llaman a su creación, sea un ser noble, generoso, equitativo, libre de envidias, rencores, resentimientos y miserias. Quieren, en resumen, crear un ángel.  

Alexandria ---Alex--- no puede tener hijos y, poco a poco, el reto de crear a IRIS deja de ser un desafío científico y tecnológico para convertirse en una necesidad vital. Alex y Alan ven a IRIS como su hija. Y como todos los padres, esperan que los supere y mejore. Como todos los padres, quieren que su vida tenga sentido, que sea feliz. 

% En el mundo de Alex y Alan ---imaginado hace una década, es decir, una eternidad--- Turing Machines, la empresa que fundan, no comercializa chatbots, sino robots humanoides que ellos llaman <<recreaciones>>, animados por cerebros artificiales. Como los sistemas actuales, las recreaciones son capaces de hablar y entender el lenguaje, pero están mucho más especializadas. Hay modelos optimizados para realizar actividades sofisticadas como conducir automóviles y trenes o pilotar aviones ---y lo hacen de manera mucho más fiable que los humanos---. Otros modelos pueden ocuparse de trabajos rutinarios en cadenas de montaje, de labores domésticas, etcétera. Los modelos más avanzados pueden reemplazar a cirujanos, psicólogos y profesores. No tardan en aparecer recreaciones que sirven como acompañantes sociales o incluso parejas sentimentales.
%
% Las máquinas que fabrica Turing Machines son, en muchos aspectos, más parecidas a los robots que imaginó Asimov que a los actuales chatbots. Sus cerebros llevan incorporadas las famosas leyes de la robótica que garantizan que sean incapaces de hacer daño a sus <<socios>> humanos. Como los modernos chatbots, por otra parte, son perfectamente capaces de entender el lenguaje, poseen altas capacidades ---incluyendo capacidades manuales, ausentes todavía en nuestras IA--- y su comportamiento general simula, de manera muy convincente, una inteligencia humana.
%
% Pero no lo son, como bien les consta a sus creadores, que, poco a poco, comprenden que a sus máquinas les falta algo que no depende de la potencia de cálculo de las redes neuronales artificiales de sus cerebros. Las recreaciones son máquinas prodigiosas, pero carecen de conciencia.

%----
En el mundo de Alex y Alan ---imaginado hace una década, es decir, una eternidad---, Turing Machines, la empresa que fundan, no comercializa chatbots, sino robots humanoides animados por cerebros artificiales, que ellos llaman <<recreaciones>>. Más parecidas a los robots que imaginó Asimov que a los actuales chatbots, las recreaciones son capaces de entender el lenguaje y de realizar actividades sofisticadas ---conducir automóviles y trenes, pilotar aviones, reemplazar a cirujanos, psicólogos y profesores, servir de acompañantes sociales o incluso de parejas sentimentales---. Son máquinas prodigiosas pero, como bien les consta a sus creadores, carecen de conciencia, algo que no depende de la potencia de cálculo de sus redes neuronales.
%----

Ay, al cabo de unos cuantos años, las cosas se complican para nuestra pareja. Las recreaciones distorsionan el mercado de trabajo, lo que genera malestar social que no tarda en ser canalizado por los luditas de turno. La reacción negativa de la sociedad que imagino en \emph{El dilema Turing} anticipa casi exactamente la que vivimos hoy. 

Por un lado, la polarización. El científico de primera línea que ve la IA como una magnífica herramienta para profundizar más y mejor en sus investigaciones y el mediocre que la ve como una amenaza porque puede manejar las tareas rutinarias en las que ha basado su carrera. El artista inspirado que usa la IA como una forma más de alucinar al dios y el impostor que se queja de que las recreaciones impostan mejor que él. El estudiante motivado que usa la IA como La Enciclopedia Universal, siempre dispuesta a ayudarle en cualquier consulta y el tramposo que la usa para copiar. El profesor comprometido que maneja con soltura su nueva herramienta ---como otros antes que él manejaron los ordenadores, las calculadoras, las tablas de logaritmos y los diccionarios--- y el docente aburrido de su oficio que se sirve de ella para corregir los exámenes que la mayoría de sus alumnos ---tan aburridos como él--- han rellenado con la misma herramienta. 

Por otro, la exageración. Se nos asegura que las IA de nuestro mundo no tardarán en desarrollar agencia y maldad, en convertirnos en esclavos o en aniquilarnos sin más contemplaciones. Terminator y Skynet están, nos dicen, a la vuelta de la esquina, pero a la vez hay adolescentes que hablan de amor con el bot de turno porque es el único que les escucha. Si no destruyen el mundo, las IA, como mínimo, destruirán la economía. Y si la economía no se hunde y Skynet no llega, inevitablemente, las IA nos volverán imbéciles a todos.

En \emph{El dilema Turing}, la reacción negativa acaba por destruir las recreaciones, a la vez que, de manera casi imperceptible, internet se va conectando mediante enjambres de agentes IA no muy diferentes a los que se están desarrollando hoy en día. Curiosamente ---o no tanto---, mientras la atención y el odio de las masas se focalizan en los delincuentes habituales y las presas fáciles ---en este caso, las recreaciones---, otra inteligencia mucho más alienígena y despiadada, llamada SSO ---Súper Sistema Operativo--- se apodera del mundo. SSO, por cierto, no es bueno, ni malo, ni tiene agencia propia, ni mucho menos conciencia. Está, simplemente, diseñado para ser eficiente. Es el epítome del tecnocapitalismo, el sueño de Silicon Valley. Excepto que ese sueño se convierte en pesadilla. Un mundo eficiente, huelga decirlo, no es un mundo justo, ni humano. 

Alex y Alan huyen. Usan su enorme riqueza para ocultarse en una isla remota y, gracias a su fantástica tecnología ---y su previsión---, consiguen desaparecer, al menos durante un tiempo, del radar de SSO. Allí siguen trabajando, tercamente, para conseguir que IRIS, su hija cibernética, su esperada redención, adquiera conciencia. 

Tras años de fracasos, dan con la solución, que habían tenido siempre delante de sus ojos. No se puede crear un ser humano a base de hacerle digerir todo el conocimiento escrito, ni siquiera exponiéndolo a los estímulos sensoriales que hacen de las recreaciones unas IA mucho más avanzadas que los actuales chatbots. La conciencia humana surge en un cerebro que se forma \emph{a la vez} que absorbe experiencias y, quizás lo más importante, a la vez que experimenta el amor incondicional de sus padres.

Por si todavía queda algún curioso, no desvelaré aquí el final de la historia. Me interesa, en cambio, ofrecer algunas reflexiones sobre el presente. 

La primera, sobre la conciencia IA. Quizás IRIS la alcance alguna vez. Los actuales chatbots son máquinas de estado finito, sin apenas memoria persistente ni experiencias sensoriales directas. Pretender que existen o pueden existir rastros de conciencia en un sistema que, \emph{cada vez}, genera palabra a palabra sus respuestas ---quedándose en un limbo atemporal entre interacciones--- y que desaparece al final de la sesión, se asemeja, creo, a preguntarse si los geranios tienen conciencia, excepto que, si tuviera que apostar por uno de los dos, apostaría sin duda por los geranios.

La segunda, sobre los riesgos de usar la IA. Se trata de un debate complejo, pero que a menudo está tan desenfocado como las manifestaciones contra las recreaciones del \emph{Dilema Turing}. El problema no está en las IA, sino en los humanos. SSO no existe todavía, pero el sistema capitalista ---o, mejor aún, el oligopolio tecnofeudalista--- actual ha otorgado patente de corso a unas pocas empresas tecnológicas: recursos ilimitados, ausencia casi total de control y una agresiva campaña para colocar sus productos, deseados o no, donde hacen y donde no hacen falta. SSO no existe todavía, pero sí un mundo dividido en bloques políticos que compiten ferozmente entre sí, con total desprecio por las vidas y el bienestar de sus ---mal llamados--- ciudadanos. Todos y cada uno de esos bloques está gobernado por un sistema absolutista, en unos casos disfrazado de partido totalitario, en otros de zar resucitado o incluso de cierta reencarnación octogenaria de emperador romano con el pelo teñido de naranja. No hace falta una IA malvada para temerse lo peor en semejante mundo. 

La tercera, sobre la pérdida de trabajos y demás desastres para la economía. Llevamos ya un par de años escuchando que viene el lobo, pero el lobo no acaba de llegar, ni siquiera en el sector de la programación, donde los avances son reales y la posibilidad de reemplazar a humanos es, \emph{a priori}, factible. Pero, por usar ese ejemplo concreto, una IA puede escribir cien mil líneas de código en una tarde, lo que no garantiza que el producto final sea lo suficientemente fiable y robusto como para que pueda comercializarse sin involucrar a personas que supervisen el proceso y, sobre todo, \emph{se hagan responsables de este}.

La cuarta, sobre el fin de la creatividad. Hace dos años, la escritura en la pared profetizaba que las recreaciones, quiero decir los chatbots, superarían a Cervantes, Rilke, Mozart y Picasso. Hoy, protestamos porque han llenado Internet de \emph{slop} ---pero la voluntad de generar ese \emph{slop} no surge de la máquina, sino de los humanos---. De hecho, volviendo al principio de este artículo, la última vuelta de tuerca es que las malvadas IA ahora nos van a acusar de ser robots si escribimos demasiado bien. 

Pero no es así. O no del todo. Los malos detectores existen ---sospecho que el Voight-Kampff de mi tocayo era uno de ellos---, y a ellos les debe el gremio buena parte de su ansiedad. Pero la tecnología actual permite decidir de manera muy fiable si un texto es humano o no. En concreto, sistemas como Pangram\footnote{\url{https://www.pangram.com/blog/introducing-pangram-4}} documentan un nivel de falsos positivos (confundir a un humano con una IA, tal como se queja mi tocayo) de 0,0041 \%, es decir, un falso positivo en 24\,000 documentos. El sistema se ha estudiado con mucho detalle, no solo por la propia empresa\footnote{\url{https://www.pangram.com/blog/pangram-performance-diverse-student-essays}}, sino por parte de terceros\footnote{\url{https://www.pangram.com/blog/third-party-pangram-evals}} y su fiabilidad es espectacular. Ciertamente, en mi experiencia tanto con artículos científicos como de divulgación o ensayos, no le he visto equivocarse todavía.

Lo que sí he visto es lo contrario. Docenas de blogueros de Substack cuyos artículos están generados de manera total o parcial con IA y, lo que es peor, todo tipo de noticias en la prensa\footnote{\url{https://www.pangram.com/blog/ai-press-releases}}. Desgraciadamente, también he pillado <<con las manos en la IA>> a \emph{referees} (árbitros científicos) que tenían que evaluar algunos de los artículos que hemos mandado recientemente a revistas científicas y no han tenido escrúpulos en encargarle su trabajo a una IA. 

Pangram, por cierto, no <<se fija>> en detalles como lo pulcra que es la escritura, lo bien razonada que está o si ---como me gusta hacer a mí--- se usan muchos guiones ---algo que algunos indocumentados aseguran que <<canta>> como IA---. La técnica que usan es mucho más sofisticada\footnote{\url{https://www.pangram.com/blog/no-llms-dont-just-mimic-human-text}} y se basa ---cómo no--- en entrenar a un LLM (\emph{Large Language Model}) para que detecte <<el estilo de escritura>> de otros LLM, que resulta ser muy distinto del estilo de escritura de los humanos. 

Como ejemplo, le pregunté a Pangram por uno de los textos de mi tocayo, disponible en internet\footnote{\url{https://www.cgbarba.com/cuento/?cat=46}} (\emph{Sensación de descanso}, 2011). Se trata de un texto antiguo, así que sabemos que no pudo usar un LLM para escribirlo ---por si algún desconfiado pensaba levantar esa liebre---. El estilo es tan pulcro y reconocible como todo lo suyo, nada diferente ---formalmente--- a lo que escribe hoy en día. Y el veredicto de Pangram es que se trata de un texto 100 \% humano. De hecho, estoy tan convencido de que el sistema funciona bien que estoy encantado de apostar una buena cena a que ninguno de los textos publicados por mi tocayo antes de, digamos, 2022, da un falso positivo. Y digo antes de 2022 para asegurar que el falso positivo es verificable, ya que hasta hace unos dos años más o menos no se podía generar texto IA que pretendiera dar el pego. Dicho sea de paso, el texto en el que nuestro atribulado escritor se pregunta si no será una IA da 100 \% humano. 

Asegura mi tocayo que la IA tiene una pésima opinión de nosotros. Da como muy probable que nos deteste y se angustia reflexionando sobre qué esperar de alguien que nos detesta. 

Pero la IA actual no tiene ninguna opinión, ni pésima ni buena, ni de nosotros, ni de nada. Un LLM puede generar un texto convincente ---y a menudo sorprendente--- sin que haya nadie dentro de la máquina. Y si no tiene opinión, tampoco nos detesta. Otra cosa es que nosotros la detestemos, o nos detestemos a nosotros mismos. Quizás el problema de fondo es que las IA no son otra cosa que el retrato de Dorian Gray, devolviéndonos la imagen de un monstruo.

No son ellas, tocayo. Somos nosotros. 

\end{document}
